% ===========================================================================
%  00 — Abstract + Zusammenfassung
%
%  Evidence: paper_context/claims_registry.md (verified claims only)
%  Terminology: paper_context/terminology.md
% ===========================================================================

\chapter*{Abstract}

Procedural learning in high-fidelity simulation requires learners to execute ordered,
state-dependent sequences while verifying intermediate system states, yet existing
training aids either break immersion or lack access to live simulator state. This
thesis addresses the gap by designing a telemetry-grounded tutoring runtime that
integrates procedure-pack-driven state tracking, retrieval-augmented help generation,
and a fine-tuned vision-language model (VLM) for structured visual fact extraction
from cockpit displays. A central finding is that visual fact ontology design governs
model reliability: evolving from an 8-fact to a 13-fact ontology by decomposing
compound targets into locally verifiable visual atoms reduced critical false positives
by 64--89\% on independent holdout sets. A Qwen3.5-9B LoRA adapter trained on 928
bilingual rows achieves 0.99 fact accuracy on two holdout sets with verified zero
training overlap; the identical data recipe transfers to Gemma-4-31B with consistent gains,
suggesting the approach is not tied to a single backbone. Controlled human-subject evaluation in VR
remains a planned next phase; the current work delivers the technical baseline,
benchmark artifacts, and evaluation infrastructure to support it.

\newpage
\chapter*{Zusammenfassung}

\begin{otherlanguage}{ngerman}
Prozedurales Lernen in hochrealistischer Simulation erfordert das Ausf\"uhren
geordneter, zustandsabh\"angiger Handlungsabfolgen bei gleichzeitiger
\"Uberpr\"ufung von Systemzust\"anden. Herk\"ommliche Trainingshilfen
unterbrechen jedoch entweder die Immersion oder haben keinen Zugriff auf den
Simulatorzustand. Diese Arbeit begegnet dieser L\"ucke durch eine
telemetriegest\"utzte Tutoring-Laufzeitumgebung, die paketbasierte
Zustandsverfolgung, Retrieval-Augmented-LLM-Hilfegenerierung und ein
feinabgestimmtes Vision-Language-Model (VLM) zur strukturierten Extraktion
visueller Fakten aus Cockpit-Anzeigen kombiniert. Ein zentrales Ergebnis ist,
dass das Design der visuellen Fakten-Ontologie die Modellzuverl\"assigkeit
bestimmt: Die Weiterentwicklung von einer 8-Fact- zu einer 13-Fact-Ontologie
durch Zerlegung zusammengesetzter Ziele in lokal \"uberpr\"ufbare visuelle
Atome reduzierte kritische False-Positives um 64--89\% auf unabh\"angigen
Holdout-Sets. Ein Qwen3.5-9B LoRA-Adapter, trainiert mit 928 bilingualen
Datenzeilen, erreicht eine Fact-Genauigkeit von 0,99 auf zwei Holdout-Sets
mit verifizierter Null-\"Uberlappung zum Training; dieselbe Datenrezeptur
\"ubertr\"agt sich mit konsistenten Verbesserungen auf Gemma-4-31B, was
darauf hindeutet, dass der Ansatz nicht an ein einzelnes Backbone gebunden
ist. Die kontrollierte Human-Subject-Evaluation in VR
bleibt eine geplante n\"achste Phase; die vorliegende Arbeit liefert die
technische Basis, Benchmark-Artefakte und Evaluationsinfrastruktur daf\"ur.
\end{otherlanguage}
